\documentclass[11pt]{article}
\usepackage[T1]{fontenc}
\usepackage{lmodern}
\usepackage{amsmath,amssymb,amsthm}
\usepackage{booktabs}
\usepackage[margin=1in]{geometry}
\usepackage{microtype}
\usepackage[colorlinks=true,linkcolor=blue,citecolor=blue,urlcolor=blue]{hyperref}
\hypersetup{
  pdftitle={Counting half-integral polygons with collinear interior lattice points},
  pdfauthor={Kaiwen Jin},
  pdfsubject={Computer-assisted exact enumeration of half-integral polygons}
}
\newtheorem{theorem}{Theorem}[section]
\newtheorem{proposition}[theorem]{Proposition}
\newtheorem{lemma}[theorem]{Lemma}
\theoremstyle{remark}
\newtheorem{remark}[theorem]{Remark}
\numberwithin{equation}{section}
\newcommand{\Z}{\mathbb Z}
\newcommand{\R}{\mathbb R}
\newcommand{\GL}{\operatorname{GL}}
\newcommand{\Fix}{\operatorname{Fix}}
\title{Counting half-integral polygons with collinear interior lattice points}
\author{Kaiwen Jin\\\small\href{mailto:yui.ui.siki@gmail.com}{yui.ui.siki@gmail.com}}
\date{October 14, 2026}

\begin{document}
\maketitle
\begin{abstract}
We count convex half-integral polygons with a prescribed number of collinear
interior lattice points, up to affine unimodular equivalence. For every
$i\geq3$, the count is the degree-seven polynomial conjectured by Bohnert and
Springer. The proof is computer-assisted: their strip theorem reduces the
polygons to two boundary chains on five integer rows, and a unique integer
shear normalization gives finitely many integer linear systems.
Complete rational generating functions for these systems, together with the
fixed classes of the remaining reflections, give the formula by Burnside's
lemma. Integral polygons are included, and translations are integral.
\end{abstract}

\section{Introduction}\label{sec:introduction}

A half-integral polygon is a two-dimensional convex polygon in $\R^2$ whose
vertices lie in $\tfrac12\Z^2$. This includes lattice polygons.
We identify $P$ and $P'$ when
\[
 P'=UP+t,\qquad U\in\GL_2(\Z),\quad t\in\Z^2.
\]
For an integer $i\geq3$, let $N(i)$ denote the number of these equivalence
classes having exactly $i$ interior points of $\Z^2$, all on one line.
Prescribing the denominator and the interior lattice points gives a concrete
classification problem for rational polygons. In the collinear case, the
line of interior points also supplies coordinates suited to counting.

Bohnert and Springer \cite{BS} classify rational polygons with small
denominator and few interior lattice points. Their Proposition~3.5 puts a
$k$-rational polygon with collinear interior lattice points in
$\R\times[-1,1]$ when the number of those points exceeds $k$.
Here $k$-rational means that multiplication by $k$ makes every vertex
integral; the denominator need not be exactly $k$.
Their Classification~3.6 records the half-integral counts for
$3\leq i\leq20$, and their Conjecture~3.7 gives the polynomial below.
We give a computer-assisted proof over the whole stated range.

\begin{theorem}\label{thm:main}
For every integer $i\geq3$,
\begin{equation}\label{eq:main}
\begin{split}
 N(i)=\frac{i+1}{1260}\bigl(&512i^6+12928i^5+137740i^4+685145i^3\\
                         &+1582743i^2+1665222i+710640\bigr).
\end{split}
\end{equation}
Here $N(i)$ counts two-dimensional convex polygons with vertices in
$\tfrac12\Z^2$, including integral polygons, with exactly $i$ collinear
interior points of $\Z^2$, modulo $\GL_2(\Z)$ and translations in $\Z^2$.
\end{theorem}

The strip theorem leaves only five possible vertex rows after doubling.
It does not by itself remove the infinitely many integer shears.
Normalizing one boundary slope removes those shears uniquely and leaves
finitely many coordinate choices at each $i$.
The remaining equivalences form a four-element group. Its stabilizers matter,
so counting normalized polygons and simply dividing by four would be
incorrect. We instead compute the generating functions of the four fixed
sets. Each is a finite sum of integer-point generating functions, with the
number of interior points retained as an unbounded variable.

The rational-series calculation uses Normaliz through PyNormaliz
\cite{BIS,Normaliz,PyNormaliz}. The proof therefore relies on trusted exact
software for a finite collection of complete generating functions.
The geometric correspondence, the fixed-point equations and the final
coefficient identity are given below. The accompanying computational
supplement contains the exact systems, complete rational outputs and
reproduction commands.

Section~\ref{sec:chains} gives the boundary description and shear
normalization. Section~\ref{sec:group} determines the remaining equivalences
and fixed sets. Section~\ref{sec:series} performs the generating-function
count and recovers~\eqref{eq:main}. Throughout, lattice points mean points of
$\Z^2$; only integer translations of the original polygon are allowed.

\section{Normalized boundary chains}\label{sec:chains}

\subsection{The interior row and the doubled polygon}

We use the following published result with its original parameter range.

\begin{proposition}[Bohnert--Springer, Proposition 3.5]\label{prop:strip}
A two-dimensional convex $k$-rational polygon with collinear interior
lattice points and more than $k$ such points is affinely unimodularly
equivalent to a polygon contained in $\R\times[-1,1]$.
\end{proposition}

For its proof see \cite[pp.~1248--1249]{BS}.
Apply Proposition~\ref{prop:strip} with $k=2$.
The only integer horizontal row strictly inside the strip is $y=0$.
The interior lattice points on that row are consecutive: convexity puts every
point between the first and last in the interior. After an integer horizontal
translation, they are
\[
 (j,0),\qquad 1\leq j\leq i.
\]
Set $n=i+1\geq4$ and $Q=2P$. Then $Q$ has integral vertices and its interior
points of the even lattice $2\Z^2$ are exactly
\begin{equation}\label{eq:interior-row}
 (2j,0),\qquad 1\leq j\leq n-1.
\end{equation}
The lowest and highest vertex heights of $Q$ are $-a$ and $b$, with
$a,b\in\{1,2\}$, since the interior row lies strictly between them.

Write the horizontal sections as
\begin{equation}\label{eq:section}
 Q\cap(\R\times\{y\})=[f(y),\,2n-g(y)]\times\{y\},
 \qquad -a\leq y\leq b.
\end{equation}
The left boundary coordinate $f$ and the right boundary distance $g$ are
convex piecewise affine functions. Their breakpoints have integer heights
and integer values. Each boundary's height set contains $-a,b$ and any
subset of the intermediate integer heights. A height is included in this
set precisely when the boundary has a strict change of slope there.
The left and right height sets can differ. At an extreme height the
horizontal face may have length zero.

\subsection{Exact integer conditions}

The row condition~\eqref{eq:interior-row} is equivalent to
\begin{equation}\label{eq:intercepts}
 0\leq f(0)<2,\qquad 0\leq g(0)<2.
\end{equation}
Indeed, the interior row section must contain $2,4,\ldots,2n-2$ strictly and
must not contain $0$ or $2n$ strictly.
For three successive boundary heights $u<v<w$, write the boundary values
as $x_u,x_v,x_w$. Strict slope increase is equivalent to the integer
inequality
\begin{equation}\label{eq:convex}
 (w-v)x_u-(w-u)x_v+(v-u)x_w\geq1.
\end{equation}
For the unique consecutive heights $u\leq0<v$ containing the right side of
zero, put $d=v-u$. Linear interpolation gives
\[
 x(0)=\frac{v x_u-u x_v}{d}.
\]
Thus the corresponding intercept condition is
\begin{equation}\label{eq:integer-intercept}
 0\leq v x_u-u x_v\leq2d-1.
\end{equation}
This applies separately to each boundary, with its own central edge.
The boundaries do not cross if and only if
\begin{equation}\label{eq:width-endpoints}
 f(-a)+g(-a)\leq2n,\qquad f(b)+g(b)\leq2n.
\end{equation}
To see sufficiency, $f+g$ is convex, hence is no larger than the larger of its
endpoint values.

\begin{lemma}\label{lem:encoding}
For fixed $n\geq4$, conditions~\eqref{eq:convex},
\eqref{eq:integer-intercept} and~\eqref{eq:width-endpoints}, with integer
boundary values and the stated height sets, describe exactly the normalized
polygons $Q$ satisfying~\eqref{eq:interior-row}. The description is unique.
\end{lemma}

\begin{proof}
For admitted chain data let
\[
 K=\{(x,y):-a\leq y\leq b,\ f(y)\leq x\leq2n-g(y)\}.
\]
This is convex, being the intersection of the epigraph of $f$ with the
hypograph of $2n-g$ and the height strip.
Its width $w(y)=2n-f(y)-g(y)$ is nonnegative by
\eqref{eq:width-endpoints} and is concave. At zero,
$w(0)>2n-4\geq4$.
Interpolating between zero and either endpoint shows that $w$ is positive
at every nonendpoint height. Hence $K$ is two-dimensional.
Its vertices are the strict side breakpoints and the endpoints of the
extreme faces, all integral. A point extreme face contributes one vertex,
although the point belongs to both chains.

Condition~\eqref{eq:intercepts} gives exactly the interior even-lattice
points in~\eqref{eq:interior-row}. The only other possible even rows are
$y=-2,2$, and each is absent or is an extreme boundary row.
Thus $K/2$ is an admissible half-integral polygon.
Conversely, every normalized polygon has these boundaries and satisfies
the conditions. Strict slope increase excludes redundant collinear
vertices, so its height sets and coordinates are recoverable uniquely.
This also retains corners at height zero and zero-length extreme faces.
\end{proof}

\subsection{Removing integer shears}

An integer shear is
\[
 S_s(x,y)=(x+s y,y),\qquad s\in\Z.
\]
It sends $(f,g)$ to $(f+s y,g-s y)$.
We choose the unique representative with
\begin{equation}\label{eq:gauge}
 0\leq f'_+(0)<1.
\end{equation}
Adding an integer to the right derivative puts it in this half-open
interval in exactly one way.
If $(u,v)$ is the left central edge, condition~\eqref{eq:gauge} says
\[
 0\leq x_v-x_u\leq d-1,\qquad d=v-u.
\]
For $\delta=x_v-x_u$ the reference values satisfy
\begin{equation}\label{eq:gauge-pairs}
 \delta\in\{0,\ldots,d-1\},\qquad
 0\leq d x_u-u\delta<2d,\qquad x_v=x_u+\delta.
\end{equation}
For each $\delta$, the half-open interval for $x_u$ has length two and
contains exactly two integers. There are therefore $2d$ reference pairs.
Also
$x_u=f(0)+(u/d)\delta\in(-2,2)$, so the pairs can be found by checking
$x_u=-2,-1,0,1,2$ and retaining exactly~\eqref{eq:gauge-pairs}.

\begin{lemma}\label{lem:finite}
For every $n\geq4$ there are finitely many normalized chain data satisfying
\eqref{eq:gauge}. The integer shear action on the polygons is free.
\end{lemma}

\begin{proof}
The right derivative of a convex piecewise affine function is a supporting
slope at zero, even when zero is a corner. Consequently
\[
 f(y)\geq f(0)+f'_+(0)y\geq-2.
\]
The endpoint widths give $g(-a),g(b)\leq2n+2$.
Convexity gives the same upper bound throughout.
For $y<0$, interpolation of $g$ between $y$ and $b$ yields
\[
 g(y)\geq \frac{b-y}{b}g(0)+\frac{y}{b}g(b)\geq-4n-4.
\]
For $y>0$, interpolation between $-a$ and $y$ gives
\[
 g(y)\geq \frac{a+y}{a}g(0)-\frac{y}{a}g(-a)\geq-4n-4.
\]
The endpoint widths now bound the endpoint values of $f$ above by
$6n+4$, and convexity gives this upper bound throughout.
All boundary coordinates thus lie in bounded intervals depending
linearly on $n$. There are finitely many height sets, so the integer
fiber at $n$ is finite. The same estimates bound the real fibers.

If $S_sQ=Q$ with $s\ne0$, choose a point of $Q$ at nonzero height.
Its iterated shear images have unbounded horizontal coordinates, contrary
to boundedness of $Q$. Hence no nonzero shear fixes $Q$.
\end{proof}

\section{Equivalences and fixed classes}\label{sec:group}

\subsection{The residual four-element group}

An equivalence between two normalized original polygons must preserve the
set $\{(j,0):1\leq j\leq n-1\}$.
This set determines its line, so the linear part of the equivalence has the
form
\[
 \begin{pmatrix}\varepsilon&s\\0&\eta\end{pmatrix},
 \qquad \varepsilon,\eta\in\{1,-1\},\quad s\in\Z.
\]
Preserving the first and last interior points forces the translation to be
$(0,0)$ if $\varepsilon=1$ and $(n,0)$ if $\varepsilon=-1$.
In doubled coordinates the latter translation is $(2n,0)$.
Thus all normalized equivalences are generated by the shears and
\begin{equation}\label{eq:reflections}
 h(x,y)=(2n-x,y),\qquad v(x,y)=(x,-y).
\end{equation}
Both $h$ and $v$ conjugate $S_s$ to $S_{-s}$, and $hv$ commutes with all
shears. Quotienting by the normal shear subgroup leaves the Klein four group
$G=\{e,h,v,hv\}$.

Let $X_n$ be the finite set of shear classes at $n$, represented by
\eqref{eq:gauge}. The original polygon classes are exactly the $G$-orbits
in $X_n$. Every action preserves $n$, including the affine map $h$.

\subsection{Fixed shear classes}

If a shear class is fixed by $\sigma\in G$, then
$\sigma Q=S_mQ$ for an integer $m$.
By freeness in Lemma~\ref{lem:finite}, this integer is unique.
Comparing the two boundary functions gives the following conditions.

For $h$, the left and right height sets must agree and
\begin{equation}\label{eq:h-fixed}
 g(y)=f(y)+m y.
\end{equation}
Indeed, $hQ$ has left coordinate $g(y)$, whereas $S_mQ$ has left
coordinate $f(y)+m y$; the right-coordinate equality is the same equation.

For $v$, one must have $a=b$, each height set invariant under negation, and
\begin{equation}\label{eq:v-fixed}
 f(-y)=f(y)+m y,\qquad g(-y)=g(y)-m y.
\end{equation}
The equality $vQ=S_mQ$ compares sections at $y$ and $-y$.
These equations need only be imposed at the paired vertex heights:
both sides are affine on the same intervening intervals.

For $hv$, one must have $a=b$ and the right height set equal to the negative
of the left one. If $hvQ=S_mQ$, applying $hv$ twice and using commutation
with shears gives $Q=S_{2m}Q$. Freeness forces $m=0$.
The remaining condition is
\begin{equation}\label{eq:hv-fixed}
 g(-y)=f(y).
\end{equation}
Conversely, these respective conditions give $hQ=S_mQ$, $vQ=S_mQ$,
or $hvQ=Q$, and hence a fixed shear class.
We impose them on the representatives with the unique left slope
normalization; they do not require the
polygon itself to be fixed before taking the shear quotient.

For each $\sigma\in G$ define the formal generating function
\[
 F_\sigma(z)=\sum_{n\geq4}|\Fix_{X_n}(\sigma)|z^n.
\]
These coefficients are finite by Lemma~\ref{lem:finite}.
Burnside's formula gives
\begin{equation}\label{eq:burnside}
 C(z):=\sum_{n\geq4}N(n-1)z^n
       =\frac{F_e(z)+F_h(z)+F_v(z)+F_{hv}(z)}4.
\end{equation}
For completeness, count pairs $(\sigma,x)$ with $\sigma x=x$.
An orbit with stabilizer size $s$ has $4/s$ elements and contributes
$(4/s)s=4$ pairs. Dividing the total by four counts the orbits, including
those with nontrivial stabilizers.

\section{Exact generating functions and the count}\label{sec:series}

\subsection{The finite systems}

The boundary description gives explicit integer linear systems with $n$
as an unbounded integer variable. For $e$, enumerate $(a,b)\in\{1,2\}^2$,
both height subsets containing the endpoints, and every left reference pair
in~\eqref{eq:gauge-pairs}. Fix those two left values; keep the other side
coordinates and $n$ as variables. Impose $n\geq4$, all strict slope
inequalities~\eqref{eq:convex}, both intercept
inequalities~\eqref{eq:integer-intercept}, and the endpoint
widths~\eqref{eq:width-endpoints}.

For $h$, require equal height sets and substitute~\eqref{eq:h-fixed},
keeping its integer $m$ as a variable.
For $v$, require $a=b$ and the two symmetric height sets, and add
\eqref{eq:v-fixed} as integer equations.
For $hv$, require $a=b$, reflect the left height set for the right one,
and substitute~\eqref{eq:hv-fixed}.
All cases retain the same left reference-pair rule and all boundary
conditions. The integer $m$, when present, is unique and bounded at each
$n$ by the boundary coordinates at a nonzero height.

These loops give the disjoint cells in Table~\ref{tab:cells}.
Disjointness follows from the unique height sets, reference pair and
shear representative. In a fixed set it also uses uniqueness of $m$.
Every integer point of a cell gives one of the appropriate fixed classes,
and every fixed class belongs to exactly one cell.

\begin{table}[ht]
\centering
\begin{tabular}{crrrr}
\toprule
$(a,b)$ & $e$ & $h$ & $v$ & $hv$\\
\midrule
$(1,1)$ & 12 & 6 & 12 & 6\\
$(1,2)$ & 64 & 16 & 0 & 0\\
$(2,1)$ & 56 & 14 & 0 & 0\\
$(2,2)$ & 288 & 36 & 72 & 36\\
\midrule
Total & 420 & 72 & 84 & 42\\
\bottomrule
\end{tabular}
\caption{Number of integer linear systems in each fixed-class count.
These are input cells, not numbers of polygons. Empty cells contribute zero.}
\label{tab:cells}
\end{table}

The supplement's \texttt{src/chain\_models.py} implements precisely these
loops. A row $[a_1,\ldots,a_r,c]$ encodes
$\sum_j a_j x_j+c\geq0$ or $=0$, with every $x_j\in\Z$.
The grading is the unit vector selecting the original $n$ coordinate.
No coordinate-sign or parity restriction is added.
There are 618 cell members but 608 literally distinct complete matrix and
grading inputs. Equal complete inputs share a computed rational function;
every original cell retains its multiplicity in the sum.

\subsection{The all-parameter series calculation}

For a cell with inequalities $A x+c\geq0$ and equations $B x+d=0$,
homogenization replaces the constants by $c\lambda,d\lambda$ and adds
$\lambda\geq0$. Its integer points at $\lambda=1$ are exactly the original
integer solutions. The $n$ grading extends by zero on the homogenizing
coordinate. The real degree fibers are bounded by the estimates in
Lemma~\ref{lem:finite}; the fixed-class equations do not change this
finiteness. Thus the inhomogeneous integer-point series is well defined.

Normaliz computes these rational series by exact rational-cone enumeration
\cite[Section~5.3]{BIS}. A disjoint semi-open simplicial decomposition and the
integer points in each fundamental domain give monomials divided by the
products of $1-z^{d_j}$ belonging to the generating rays.
Extracting the $\lambda=1$ slice gives the inhomogeneous series.
The integer-lattice, row and grading conventions used here are described in
\cite[Sections~2.1, 2.9, 4.1.18--4.1.19, 4.5.1 and 4.12--4.13]{Normaliz}.
In particular, the explicit inequalities prevent the default positive
orthant convention from restricting signed coordinates, and equations are
intersected with the full integer lattice. Cells in proper affine subspaces
and empty cells are included.

The calculation uses PyNormaliz 2.24 and the collective request
\[
 \texttt{Compute(HilbertSeries, BigInt, NoGradingDenom)}.
\]
Each returned grading is checked to be the original $n$ unit vector, with
homogenizing zero and denominator one.
\texttt{NoGradingDenom} preserves this grading, including affine lattices
with nontrivial period. The result is a complete rational function over all
$n\geq4$; no coefficient cutoff is used in its computation.

\begin{lemma}[Computer-assisted rational-series calculation]\label{lem:series}
Summing the complete integer-point series of the cells in
Table~\ref{tab:cells}, with all original multiplicities, gives
\begin{align}
 F_e(z)&=\frac{z^4}{(1-z)^8}
 \bigl(669807-3348976z+7783173z^2-10577048z^3 \notag\\
 &\hspace{38mm}+8928985z^4-4635736z^5+1361923z^6-173936z^7\bigr),
 \label{eq:Fe}\\
 F_h(z)&=\frac{z^4(2349-7519z+9303z^2-5255z^3+1138z^4)}
                    {(1-z)^5},\label{eq:Fh}\\
 F_v(z)&=\frac{z^4(2069-6465z+7915z^2-4447z^3+960z^4)}
                    {(1-z)^5},\label{eq:Fv}\\
 F_{hv}(z)&=\frac{z^4(335-806z+669z^2-190z^3)}
                    {(1-z)^4}.\label{eq:Fhv}
\end{align}
\end{lemma}

\begin{proof}[Computational verification]
The supplement records the exact input, whole rational output and original
grading for each of the 608 distinct systems, and the assignment of all
618 cell members to those systems.
All systems were computed in ordinary \texttt{BigInt} mode.
An integer-only driver sums the functions with their cell multiplicities:
the returned denominator degrees are one or two, so
$1-z^2=(1-z)(1+z)$ gives a common denominator using integer polynomial
multiplication and addition.
The four rational differences from~\eqref{eq:Fe}--\eqref{eq:Fhv} are zero.
The driver checks full numerator identities rather than a finite number of
power-series coefficients.
\end{proof}

This lemma is the proof's software-dependent step. The supplement pins the
actual wheel and native-library bytes.
According to \cite[Section~5.3 and Appendix~E.2.10]{Normaliz},
\texttt{BigInt} avoids the initial 64-bit main computation.
It does not enlarge bounded degree types, and ordinary LLL arithmetic remains
part of the implementation. The computation completed normally; its observed
output degrees do not bound every intermediate operation.
The manual version is 3.11.1, while the exact minor version of the compiled
Normaliz library in the wheel is not exposed and is not inferred from it.
The calculation relies on this trusted software and runtime; an independent
counting implementation or a separately checked simplicial-decomposition
certificate is not supplied.

\subsection{Recovering all coefficients}

Combining Lemma~\ref{lem:series} with~\eqref{eq:burnside} gives
\begin{equation}\label{eq:C}
 C(z)=\frac{z^4 P_*(z)}{(1-z)^8},
\end{equation}
where the numerator polynomial is
\begin{equation}\label{eq:numerator}
\begin{split}
 P_*(z)={}&168640-844590z+1965375z^2-2673454z^3\\
          &+2258540z^4-1173244z^5+344837z^6-44056z^7.
\end{split}
\end{equation}
This is an identity of rational functions, hence of their formal
expansions at zero.

\begin{proof}[Proof of Theorem~\ref{thm:main}]
Write $c_k$ for the coefficient of $z^k$ in~\eqref{eq:numerator}.
The formal geometric-series identity gives
\[
 (1-z)^{-8}=\sum_{r\geq0}\binom{r+7}{7}z^r.
\]
Since the degree variable is $n=i+1$, extracting $z^{i+1}$
in~\eqref{eq:C} yields
\begin{equation}\label{eq:binomial-count}
 N(i)=\sum_{k=0}^7 c_k\binom{i+4-k}{7},\qquad i\geq3.
\end{equation}
If a shifted exponent is negative, its coefficient is zero and the
corresponding binomial upper argument is in $\{0,\ldots,6\}$.
The polynomial product defining that binomial also vanishes.
The smallest upper argument is $i-3\geq0$, so this formula includes every
endpoint case in the stated range.

Multiplying~\eqref{eq:binomial-count} by $7!=5040$ and expanding gives the
integer polynomial identity
\begin{equation}\label{eq:polynomial-identity}
\begin{split}
 \sum_{k=0}^7 c_k\prod_{j=0}^6(i+4-k-j)
 =4(i+1)\bigl(&512i^6+12928i^5+137740i^4+685145i^3\\
             &+1582743i^2+1665222i+710640\bigr).
\end{split}
\end{equation}
Both sides have ascending coefficients
\[
 (2842560,\ 9503448,\ 12991860,\ 9071552,\
   3291540,\ 602672,\ 53760,\ 2048).
\]
Dividing by 5040 proves~\eqref{eq:main}, since $5040=4\cdot1260$.
The degree-seven identity follows from the complete rational series;
finite classification data are not used to infer it.
\end{proof}

\paragraph{Computational materials.}
The accompanying versioned supplement contains the chain generator,
the complete compact matrix/series data, an integer-only aggregation
verifier, a Normaliz replay driver, finite integer-geometry controls and
environment information. Its README distinguishes verification of the
stored algebra from recomputation of the software-dependent series.
The full replay is optional for using the supplied data and is required
to reproduce that series calculation from the inputs.

\paragraph{AI-use statement.}
OpenAI Codex was used to develop the boundary encoding, implement and run the
exact computations, check the argument and prepare the manuscript.
Any accompanying manuscript review identifies its AI reviewer and the
reviewed version. AI review and successful software tests do not constitute
independent human mathematical certification.

\begin{thebibliography}{9}
\bibitem{BS}
M.~Bohnert and J.~Springer,
\emph{Classifying rational polygons with small denominator and few interior
lattice points}, Discrete \& Computational Geometry \textbf{76} (2026),
1238--1265.
\href{https://doi.org/10.1007/s00454-026-00822-0}{doi:10.1007/s00454-026-00822-0}.

\bibitem{BIS}
W.~Bruns, B.~Ichim and C.~S\"oger,
\emph{The power of pyramid decomposition in Normaliz},
Journal of Symbolic Computation \textbf{74} (2016), 513--536.
\href{https://doi.org/10.1016/j.jsc.2015.09.003}{doi:10.1016/j.jsc.2015.09.003}.
The algorithm description used here is Section~5.3 of
\href{https://arxiv.org/abs/1206.1916v4}{arXiv:1206.1916v4}.

\bibitem{Normaliz}
W.~Bruns and M.~Horn,
\emph{Normaliz 3.11.1}, documentation.
\url{https://github.com/Normaliz/Normaliz/blob/v3.11.1/doc/Normaliz.pdf}.

\bibitem{PyNormaliz}
\emph{PyNormaliz}, version 2.24.
\url{https://github.com/Normaliz/PyNormaliz/tree/v2.24}.
\end{thebibliography}
\end{document}
